\documentclass[11pt,a4paper]{article}
\usepackage[utf8]{inputenc}
\usepackage[T1]{fontenc}
\usepackage[english]{babel}
\usepackage{amsmath,amssymb,amsthm}
\usepackage{amsfonts}
\usepackage{mathtools}
\usepackage{geometry}
\usepackage{hyperref}
\usepackage{enumitem}
\usepackage{booktabs}
\usepackage{array}
\usepackage{mathrsfs}
\usepackage{calrsfs}
\usepackage{xcolor}
\usepackage{microtype}
\usepackage{tikz}
\usepackage{tikz-cd}
\usetikzlibrary{matrix,arrows}

\geometry{a4paper, top=2.5cm, bottom=2.5cm, left=2.5cm, right=2.5cm}

\theoremstyle{plain}
\newtheorem{theorem}{Theorem}
\newtheorem{lemma}[theorem]{Lemma}
\newtheorem{proposition}[theorem]{Proposition}
\newtheorem{corollary}[theorem]{Corollary}
\newtheorem{conjecture}[theorem]{Conjecture}

\theoremstyle{definition}
\newtheorem{definition}[theorem]{Definition}
\newtheorem{example}[theorem]{Example}
\newtheorem{remark}[theorem]{Remark}

\DeclareMathOperator{\Ext}{Ext}
\DeclareMathOperator{\Hom}{Hom}
\DeclareMathOperator{\Ker}{Ker}
\DeclareMathOperator{\Coker}{Coker}
\DeclareMathOperator{\Spec}{Spec}
\DeclareMathOperator{\Proj}{Proj}
\DeclareMathOperator{\Gal}{Gal}
\DeclareMathOperator{\Aut}{Aut}
\DeclareMathOperator{\End}{End}

\title{\textbf{The Nagata Embedding Theorem}}
\author{Pierre Deligne}
\date{}

\begin{document}

\maketitle

\noindent\textbf{Translator's note.} This is a working English translation of the uploaded LaTeX transcription of Deligne's \emph{Le th\'eor\`eme de plongement de Nagata}. It is intended as a literal mathematical translation, not a modern recast. Mathematical notation and the structure of the transcription have been kept as close as possible to the source.

\bigskip
\noindent\textbf{Abstract.} Except for a few corrections added in 2010, this article reproduces notes sent to Nagata in the 1970s, translating his papers [2] and [3] into the language of schemes. The proofs are a ``constructive'' version of Nagata's proofs. Making the proofs constructive made it possible to avoid the use of valuations. A more detailed exposition and additional results may be found in Conrad [1]. Unless expressly stated otherwise, all schemes considered below are assumed to be Noetherian.

\section*{0. Recollections on blow-ups}

\subsection*{0.1.}

Let $\mathfrak{a}$ be an ideal of the ring $A$. The scheme obtained from $X=\Spec(A)$ by blowing up $\mathfrak{a}$ is the union of the open sets $D^+(x)$, where $x$ ranges over a system of generators of $\mathfrak{a}$, and $D^+(x)$ is the spectrum of the subring $A[\mathfrak{a}x^{-1}]$ of $A[x^{-1}]$. Moreover,
\[
D^+(x)\cap (X-V(\mathfrak{a}))=\Spec(A[x^{-1}]).
\]

Let $\mathfrak{c}$ be a second ideal of $A$, defining a closed subscheme $V(\mathfrak{c})$ of $X$. The trace on $D^+(x)$ of the proper transform of $V(\mathfrak{c})$ is defined by the ideal of $A[\mathfrak{a}x^{-1}]$ obtained as the trace of the ideal of $A[x^{-1}]$ generated by $\mathfrak{c}$.

\begin{lemma}
Let $X$ be a scheme, let $\mathfrak{a}$ and $\mathfrak{b}$ be two sheaves of ideals on $X$, and let $\widetilde{X}$ be the blow-up of $X$ along $\mathfrak{a}+\mathfrak{b}$. Then the open subset of $\widetilde{X}$ which is the union of the $D^+(x)$ for $x\in\mathfrak{b}$ contains the proper transform of $V(\mathfrak{a})$.
\end{lemma}

Indeed, $\widetilde{X}$ is the union of the $D^+(x)$ for $x\in\mathfrak{a}$ or $x\in\mathfrak{b}$, and for $x\in\mathfrak{a}$ the open set $D^+(x)$ is disjoint from the proper transform of $V(\mathfrak{a})$.

\bigskip
Lemma 0.3 of the original text was false. The corrected version that follows dates from March 1996. It led to a change of notation. The original lemma used the following notation, to which the rest of the text refers: instead of $F$ and $G$, one considered their complements $U$ and $V$, which, together with $Y' := U\cap p^{-1}(Y)$, formed a commutative diagram
\[
\begin{tikzcd}
Y' \arrow[d] \arrow[r, hook] & Z \arrow[d, "p"] \\
U \arrow[r, hook] & V \arrow[r, hook] & X .
\end{tikzcd}
\]
Here $Y'$ is closed in $p^{-1}(V)$.

\begin{lemma}
Let $p:Z\to X$ be a morphism, let $F\supset G$ be closed subschemes of $X$, and let $Y$ be a closed subscheme of $Z$. Assume that
\[
Y\cap p^{-1}(F)\subset p^{-1}(G)\qquad\text{scheme-theoretically.}
\]
Blow up $X$ along $G$ to obtain $\widetilde{X}$, and let $\tilde p:\widetilde{Z}\to\widetilde{X}$ be obtained from $p:Z\to X$ by base change. Then the closure
\[
\overline{(Y\cap p^{-1}(G)\text{ in }\widetilde{Z})}
\]
of $Y\cap p^{-1}(G)$ in $\widetilde{Z}$ and
\[
\tilde p^{-1}(\widetilde{F\cap G\text{ in }\widetilde{X}})
\]
are disjoint closed subsets of $\widetilde{Z}$.
\end{lemma}

We may assume that $Z$ and $X$ are affine. Let $G$ denote the ideal of $G$. Cover $\widetilde{X}$ by the open subsets $D^+(g)$ for $g\in G$ (a system of generators of $G$ suffices), and cover $\widetilde{X}-\widetilde{G}$ by the open sets
\[
\widetilde U(g)=D^+(g)\cap(\widetilde X-\widetilde G)\subset \widetilde X-\widetilde G,
\]
with
\[
\widetilde U(g)\subset D^+(g).
\]
It is enough to see that, for each $g$, the traces on $\tilde p^{-1}D^+(g)$ of
\[
\tilde p^{-1}(\widetilde{F\cap G\text{ in }\widetilde X})
\]
and of
\[
\overline{(Y\cap p^{-1}G\text{ in }\widetilde Z)}
\]
are disjoint. They are, respectively, the inverse image of the closure of $F\cap\widetilde U(g)$ in $D^+(g)$, and the closure of $Y\cap p^{-1}(\widetilde U(g))$ in $\tilde p^{-1}D^+(g)$.

Write $\sim$ for inverse image from $\widetilde X$ to $Z$, or from $D^+(g)$ to $\tilde p^{-1}D^+(g)$. Since $p^{-1}G\supset Y\cap p^{-1}(F)$, the inverse image $\tilde g$ of $g$ on $\widetilde Z$ can be written
\[
\tilde g=y+f,
\]
where $y$ vanishes on $Y$ and $f$ vanishes on $p^{-1}(F)$. There are functions $a_i$ on $\widetilde Z$ and elements $f_i$ in the ideal of $F$ such that
\[
f=\sum a_i f_i^{\sim}.
\]
Since $F\supset G$, the $f_i$ belong to the ideal of $G$, and the $f_i/g$ are functions on $D^+(g)$. They vanish on $F\cap\widetilde U(g)$, hence on
\[
\overline{(F\cap\widetilde U(g)\text{ in }D^+(g))},
\]
and $\sum a_i(f_i/g)^{\sim}$ vanishes on
\[
\tilde p^{-1}\overline{(\widetilde U(g)\text{ in }D^+(g))}.
\]
On $\tilde p^{-1}D^+(g)$ one has
\[
1=y/g+f/g=y/g+\sum a_i(f_i/g)^{\sim},
\]
and the same sum is therefore equal to $1$ on $Y\cap p^{-1}\widetilde U(g)$, hence on its closure
\[
\overline{(Y\cap\widetilde U(g)\text{ in }\tilde p^{-1}D^+(g))}.
\]
This proves the required disjointness.

\begin{remark}
Suppose instead that we are given closed subsets $F\supset G_0$ of $X$, with $G_0$ such that $Y\cap p^{-1}(F)$ is set-theoretically contained in $p^{-1}(G_0)$. Then there exists an $r$ such that the $r$-th infinitesimal neighborhood $G$ of $G_0$ in $F$ satisfies the hypotheses of the lemma. If $J$ is the ideal defining
\[
\overline{(Y\cap p^{-1}(F))}_{\mathrm{red}}
\]
in $Y\cap p^{-1}(F)$, it is enough to choose $r$ such that $J^{r+1}=0$.

Thus the lemma makes it possible, by a blow-up that does not affect the complement $V$ of $G_0$, to ``separate'' the proper transform of $F$ and
\[
\overline{(Y\cap p^{-1}(G_0))}
\]
in $Z$.
\end{remark}

\begin{lemma}
Consider a commutative diagram of schemes
\[
\begin{tikzcd}
\widetilde X \arrow[r,"p"] & X \\
\widetilde X' \arrow[u] \arrow[r,"q"'] & X' \arrow[u]
\end{tikzcd}
\]
with $U$ a dense open subset of both $X$ and $X'$, with $p$ and $q$ separated of finite type, and with $q^{-1}(X)=\widetilde X'$.

Let $F$ be a closed subset of $X'$, and let $G$ be a closed subset of $\widetilde X'$ containing $p^{-1}(F)$. Then there exists an ideal $\mathfrak a$ with
\[
V(\mathfrak a)\subset \overline F-F
\]
such that, after base change by the blow-up morphism defined by $\mathfrak a$, $\widetilde X\to X$, and after replacing the new $\widetilde X'$ by the closure of $U$, one has
\[
\widetilde q^{-1}(F)\subset \widetilde G.
\]
\end{lemma}

In Nagata's language, to take the closure of $F$ in the Zariski--Riemann space of $X$ (the ``bubble space''), it is enough to take the ``intersection'' of the closures of $F$ in the blow-ups of $X$ whose centers are contained in $\overline F-F$.

Apply Lemma 0.3 to
\[
\begin{tikzcd}
\widetilde X' \arrow[r] & q^*F \arrow[d] \\
& X-F \arrow[r,hook] & X-(\overline F-F) \arrow[r,hook] & X .
\end{tikzcd}
\]
This gives $\mathfrak a$. After the base change $\widetilde X\to X$, the new closure $\widetilde F$ is disjoint from the image of the closure of the old $\widetilde X'-q^*F$ in $\widetilde X'-G$. In particular, after replacing $\widetilde X'$ by $\widetilde U$, this new closure is disjoint from the image of $\widetilde X'-G$, which is what was required.

\begin{lemma}
Let there be given a diagram of separated $S$-schemes of finite type
\[
\begin{tikzcd}
& \widetilde X_i \arrow[dr] \\
U \arrow[ur,hook] \arrow[dr,hook] \arrow[r,hook] & X_i \arrow[r] & \widetilde X_i \quad (1\le i\le n) \\
& X
\end{tikzcd}
\]
with $\widetilde X_i$ proper over $S$ and $U$ a dense open subset of $X$ and of the $\widetilde X_i$. Let $X^*$ be the closure of $U$ in the product of the $\widetilde X_i$. Finally, let $F_i$ be a closed subset of $X_i$; assume that the intersection of the inverse images of the $F_i$ in $X^*$ is empty.

Then, after replacing each $X_i$ by a blow-up with center in $\overline{F_i}-F_i$, one can obtain an analogous situation in which the intersection of the inverse images of the $F_i$ is empty.
\end{lemma}

Let $X^{**}$ be a blow-up of $X^*$,
\[
X^{**}\to \prod_i \widetilde X_i,
\]
such that in $X^{**}$ one has $p_i^{-1}(F_i)=\emptyset$; this is obtained by blowing up $p_i^{-1}(F_i)$. Apply Lemma 0.4 to the diagrams
\[
\begin{tikzcd}
p_i^{-1}(\widetilde X_i) \arrow[r,hook] \arrow[d] & X^{**} \arrow[d] \\
U \arrow[r,hook] & X_i \arrow[r,hook] & \widetilde X_i
\end{tikzcd}
\]
and to $F_i\subset X_i$, $G_i=p_i^{-1}(F_i)$. This provides the desired blow-ups of the $X_i$. Indeed, in the new situation, in the closure of $U$ in $X^{**}$ one has
\[
p_i^{-1}(F_i)\subset G_i=\emptyset.
\]
Since $X^{**}\to X^*$ is surjective, the same remains true in $X^*$.

\section*{1. The theorems}

\begin{definition}
Let $X$ and $Y$ be two $S$-schemes, with $Y$ separated over $S$. A quasi-domination $f:X\dashrightarrow Y$ is a pair consisting of a dense open subset $U$ of $X$ and a morphism of $S$-schemes $f:U\to Y$ whose graph is closed in $X\times_S Y$.

Let $X$ and $Y$ be two $S$-schemes, with $Y$ separated over $S$, let $U$ be a quasi-compact open subset of $X$, and let $f:U\to Y$. We say that $X$ is quasi-dominant over $Y$ relative to $U$ and $f$ if there exists a quasi-domination of the scheme-theoretic closure $\overline U$ of $U$, with values in $Y$, extending $f$. There is at most one such quasi-domination; its graph is the closure of the graph of $f$.

A quasi-domination $f$ is called proper, or complete, if the morphism $f:U\to Y$ is proper.
\end{definition}

\begin{theorem}[{[2, Theorem 3.2]}]
Let $X$ and $Y$ be two $S$-schemes, and let $U\subset V\subset X$ be two dense open subsets of $X$. Assume that $Y$ is separated of finite type over $S$. Let $f:U\to Y$ be a quasi-domination of $V$ into $Y$. Then there exists an ideal $\mathfrak a$ with
\[
V(\mathfrak a)\subset X-V
\]
such that the blow-up $\widetilde X$ of $X$ along $V(\mathfrak a)$ quasi-dominates $Y$.
\end{theorem}

By the base change $X\to S$, one reduces to the case $X=S$. Then $Y$ is a scheme over $X$, $f$ is a section of $Y$ over $U$, and by hypothesis $f(U)$ is closed over $V$.

\smallskip
\noindent\textbf{Case 1.} $U=V$, $p:Y\to X$ quasi-affine, $X$ affine.

By hypothesis there is a commutative diagram
\[
\begin{tikzcd}
Y \arrow[rr,hook] \arrow[dr,"p"'] && E \arrow[dl] \\
& X
\end{tikzcd}
\]
with $E$ an affine space of finite type over $X$. If $X$ quasi-dominates $E$ relative to $f$, then \emph{a fortiori} $X$ quasi-dominates $Y$. Thus we may assume $Y=E_X$. A preliminary blow-up makes it possible to assume that $X-U$ is a divisor $D$. Locally on $X$, the section $f$ therefore has coordinates $(f_i u^{-n})$, where $u$ is a local equation for $D$, the $f_i$ are regular, and $n$ is chosen large enough once and for all. Let $\mathfrak a$ be the ideal generated by $u$ and the $f_i$. After blowing up $V(\mathfrak a)$, the section $f$ extends to a section of the projective completion of $E_X$, and this proves the case.

\smallskip
\noindent\textbf{Case 2.} $U=V$, $p$ quasi-affine.

By Case 1, if $(U_i)_{0\le i\le n}$ is an affine open covering of $X$, then on each $U_i$ there exists an ideal $\mathfrak a_i$ which works. This ideal has an extension $\widetilde{\mathfrak a}_i$ to $X$, with $V(\widetilde{\mathfrak a}_i)\subset X-V$, and it is enough to blow up the product of the $\widetilde{\mathfrak a}_i$.

\smallskip
\noindent\textbf{Case 3.} $p$ quasi-affine.

By Lemma 0.3, applied to
\[
\begin{tikzcd}
& f(U) \arrow[d,hook] \arrow[r,hook] & Y \arrow[d,"p"] \\
& U \arrow[r,hook] & V \arrow[r,hook] & X,
\end{tikzcd}
\]
a preliminary blow-up makes it possible to assume that
\[
p(\overline{f(U)})\cap (V-U)=\emptyset.
\]
There is therefore a quasi-compact open subset $X'$ of $X$ with $X'\cap V=U$ and $p(\overline{f(U)})\subset X'$. By Case 2, there exists an ideal $\mathfrak a$ on $X'$, with $V(\mathfrak a)\subset X'-U$, such that the blow-up $\widetilde{X'}$ quasi-dominates $Y$, equivalently $\overline{f(U)}$. It remains only to extend $\mathfrak a$ to an ideal on $X$, with $V(\mathfrak a)\subset X-V$.

\smallskip
\noindent\textbf{General case.}

Let $(Y_i)$ be a finite covering of $Y$ by affine open subsets, and set $U_i=f^{-1}(Y_i)$, $X_i=\overline U_i$, $V_i=X_i\cap V$.

Case 3 applies to the diagrams
\[
\begin{tikzcd}
U_i \arrow[r,hook] \arrow[d,"f_i"'] & V_i \arrow[r,hook] & X_i \arrow[d,hook] \\
Y_i \arrow[r,hook] & Y & X,
\end{tikzcd}
\]
and gives an ideal $\mathfrak a_i$ on $X_i$. This ideal defines an ideal, still denoted $\mathfrak a_i$, on $X$, again with $V(\mathfrak a_i)\subset X-V$. Blow up the product of the $\mathfrak a_i$. The resulting blow-up $q:\widetilde X\to X$ is such that $q^{-1}(X_i)$ dominates the blow-up $\widetilde X_i$. If $X_i'=q^{-1}(X_i)$, there exists an open subset $U_i'$ of $X_i'$, with $U_i'\supset q^{-1}(U_i)$, such that, over $X_i'$, $f(U_i')$ is a section of $f:U_i'\to Y$. We conclude that after the blow-up,
\[
f(U)=\bigcup f_i(U_i').
\]

To prove Theorem 1.2, we may replace $Y$ by $\overline{f(U)}$, and we conclude by Case 3 and the following lemma.

\begin{lemma}
$\overline{f(U)}$ is quasi-affine over $X$.
\end{lemma}

Indeed, $\overline{f(U)}$ is quasi-finite over $X$.

\begin{corollary}[A variant of Chow's lemma]
Let $f:X\to S$ be a separated morphism of finite type. Then for every open subset $U$ which is quasi-projective over $S$ and dense in $X$, there exists a diagram of $S$-schemes
\[
\begin{tikzcd}
X & \widetilde X' \arrow[l,"q"'] \arrow[r,hook,"j"] & \widetilde X
\end{tikzcd}
\]
where $q$ is a blow-up morphism defined by an ideal satisfying $V(\mathfrak a)\subset \widetilde X'-U$, and where $\widetilde X$ is projective over $S$.

In particular, $q$ is projective and surjective, and $q^{-1}(U)\xrightarrow{\sim}U$.
\end{corollary}

Let $U^*$ be a projective $S$-scheme containing $U$ as a schematically dense open subset. By Theorem 1.2 applied to the inclusion of $U$ in $X$, one can choose $U^*$ quasi-dominant over $X$: one has a diagram
\begin{equation}\label{eq:1.4.1}
\begin{tikzcd}
& U \arrow[dl,hook'] \arrow[dr,hook] \\
V \arrow[dr] && U^* \arrow[ll,"\phi"'] \arrow[dl] \\
& \overline U \arrow[r,hook] & X,
\end{tikzcd}
\end{equation}
and here, since $U^*$ is proper over $S$, $\phi$ is proper.

\begin{lemma}
Let there be given a diagram of separated $S$-schemes of finite type
\[
\begin{tikzcd}
U \arrow[d,hook] \arrow[r,hook] & V \arrow[d] \\
X & Y \arrow[l,"\phi"']
\end{tikzcd}
\]
with $U$ dense in $X$, dense in $Y$, and schematically dense in $V$, and with $\phi$ a quasi-domination of $Y$ into $X$. Let $\mathfrak a$ be an ideal of $X$, with $V(\mathfrak a)\subset X-U$, such that the blow-up $\widetilde X$ of $X$ along $\mathfrak a$ quasi-dominates $V$ relative to $U$, and such that $U$ is schematically dense in $\widetilde X$. Thus we obtain a diagram
\[
\begin{tikzcd}
& \widetilde U \arrow[dl,hook'] \arrow[dr,hook] \\
\widetilde V \arrow[dr,"\psi"'] && \widetilde W \arrow[ll,"\beta"'] \arrow[dl,"\alpha"] \\
& \widetilde X .
\end{tikzcd}
\]
Then:
\begin{enumerate}[label=(\alph*)]
\item $\psi$ is the blow-up morphism of $V$ along $\phi^*\mathfrak a$;
\item the graph of $\psi$ is closed in $\widetilde X\times_S Y$.
\end{enumerate}
\end{lemma}

Let $\widetilde V$ be the blow-up of $V$ along $\phi^*\mathfrak a$.
\begin{enumerate}[label=(\arabic*)]
\item By the universal property of $\widetilde X$, there is a unique $\widetilde X$-morphism from $\widetilde V$ to $\widetilde X$; denote it by $\alpha$.
\item By the universal property of $\widetilde V$, there is a unique $V$-morphism from $\widetilde W$ to $\widetilde V$; denote it by $\beta$.
\end{enumerate}

\[
\begin{tikzcd}
\widetilde W \arrow[r,"\beta"] \arrow[d,"\alpha"'] & \widetilde V \arrow[d,"\psi"] \\
\widetilde X & V \arrow[r,hook] & Y
\end{tikzcd}
\qquad
\begin{tikzcd}
U \arrow[r,hook] \arrow[d,hook] & \widetilde W \arrow[d] \\
V \arrow[r] & \widetilde V \arrow[r] & \widetilde X .
\end{tikzcd}
\]

By Theorem 1.2, after blowing up $X_1$ without changing $U$, we reduce to the case where $X_1$ quasi-dominates $X_2$.

\[
\begin{tikzcd}
& U \arrow[dl,hook'] \arrow[dr,hook] \\
V \arrow[dr] \arrow[rr] && X_1 \arrow[dl] \\
& X_2 .
\end{tikzcd}
\]

Apply Lemma 1.5 to this diagram so as to obtain
\[
\begin{tikzcd}
& U \arrow[dl,hook'] \arrow[dr,hook] \\
\widetilde V \arrow[dr,"\psi"'] && \widetilde{X_2} \arrow[dl] \\
& X_1,
\end{tikzcd}
\]
where
\begin{enumerate}[label=(\alph*)]
\item $\psi$ is the blow-up morphism of an ideal $\mathfrak a$ of $V$ with $V(\mathfrak a)\subset V-U$;
\item the graph of $\psi$ is closed in $X_1\times_S\widetilde{X_2}$; and
\item $\widetilde{X_2}$ is a blow-up of $X_2$.
\end{enumerate}

Let $\mathfrak a'$ be an ideal of $X_1$ extending $\mathfrak a$, with $V(\mathfrak a')\subset X_1-U$, and let $q:\widetilde{X_1}'\to X_1$ be the blow-up of $X_1$ along $\mathfrak a$. Then $\psi$ induces an isomorphism between $\widetilde W$ and $q^{-1}(V)$, and the graph of this isomorphism is closed in $\widetilde{X_1}'\times\widetilde{X_2}$. Define $\widetilde X$ by gluing $\widetilde{X_1}'$ and $\widetilde{X_2}$ along this isomorphism.

\begin{remark}
The proof shows that, for the maps $p_i:U_i\to X_i$ with $U_i\subset X_i$, one may take a blow-up morphism of an ideal supported outside $U$.
\end{remark}

\begin{theorem}[{[2, Lemma 4.1 and Theorem 4.3]}]
Let $f:X\to S$ be separated of finite type. Then there exists an immersion of $X$ into an $S$-scheme proper over $S$.
\end{theorem}

The graph of $\phi$ is, trivially, proper over $V$; the same is true for its inverse image on $V\times_S X$, and it follows that the graph of $\psi$ is proper over $V$, hence that $\psi$ is proper. Thus $\beta$ is proper, and has schematically dense image, since $U$ is already schematically dense in $V$. On the other hand, $\beta$ is a local section of $\alpha$ on $\widetilde W$ because $\beta|_U$ is one on $U$; hence $\beta$ is an isomorphism.

By hypothesis, the graph of $\psi$ is closed in $V\times_S X$, and is contained in the inverse image of the graph of $\phi$, which is closed in $Y\times X$. It follows that $\Gamma(\psi)$ is closed in $Y\times\widetilde X$.

Apply Lemma 1.5 to \eqref{eq:1.4.1}. There exists an ideal $\mathfrak a$ on $X$, with $V(\mathfrak a)\subset X-U$, such that the blow-up $\widetilde X$ quasi-dominates $\overline U\subset X-V$. Moreover, one can arrange that $U$ is schematically dense in $\widetilde X$. By Lemma 1.5, and because $\phi$ is proper, $\widetilde X$ is therefore the blow-up of $V$ along $\phi^*\mathfrak a$. If $\mathfrak a'$ is an ideal of $U^*$ extending $\mathfrak a$, with $V(\mathfrak a')\subset U^*-U$, we then take
\[
\widetilde X'=\widetilde X,\qquad \widetilde X=U^*.
\]

\begin{lemma}[{[2, Lemma 4.2]}]
Let there be given a diagram of separated $S$-schemes of finite type
\[
\begin{tikzcd}
X_2 & U \arrow[l,hook'] \arrow[r,hook] & X_1
\end{tikzcd}
\]
with $U$ a dense open subset of the $X_i$. Then there exists an embedding $j:U\hookrightarrow X$, with $X$ separated of finite type over $S$ and $U$ schematically dense in $X$, and there exist proper quasi-dominations $p_i:X\dashrightarrow X_i$.
\end{lemma}

Let $U_i$ be a finite open covering of $X$ by dense quasi-projective open subsets. For each $U_i$ there exists, by Corollary 1.4, a commutative diagram of $S$-schemes
\[
\begin{tikzcd}
& \widetilde X_i \arrow[d] \\
U_i \arrow[ur,hook,"j_i"] \arrow[r,"q_i"'] & X_i
\end{tikzcd}
\]
with $q_i$ proper, $\widetilde X_i$ proper over $S$, and $U_i$ dense in $\widetilde X_i$. Let $F_i=X_i-U_i$ and $F_i'=q_i^{-1}(F_i)$. Finally, let $X^*$ be the closure of $U$ in the product of the $\widetilde X_i$, and write $p_i:X^*\to\widetilde X_i$.

By Lemma 0.5, after a preliminary blow-up of $X_i$ which does not affect $\widetilde X_i'$, we may assume that
\[
p_i^{-1}(F_i')=\emptyset.
\]

Let $M$ be the separated $S$-scheme obtained by gluing the $\widetilde X_i$ and $X_i-F_i'$ along $U_i$.

\[
\begin{tikzcd}
X_i-F_i' \arrow[d,hook] \\
\widetilde X_i \arrow[r] & M .
\end{tikzcd}
\]

Apply Lemma 1.7 to $X$ and to the $M_i$, to obtain $X\to M$ where $M$ properly quasi-dominates each $M_i$. It remains to prove that $M$ is proper over $S$.

Indeed, one has maps
\[
X^*-p_i^{-1}(F_i')\to M_i
\]
which glue over $U=\bigcup U_i$. After replacing $X^*$ by a blow-up $X^{**}$, one obtains maps
\[
X^{**}-p_i^{-1}(F_i')\to M_i
\]
which glue. This gives a map from the scheme-theoretic closure of $U$ in $X^{**}$ with values in $M$. The scheme $M$, being the image of this map, is proper over $S$.

\section*{Comments}

I hope that, by being still more constructive, one can prove Theorem 1.6 for $f:X\to S$ separated of finite type with $S$ quasi-compact and quasi-separated. The rules to follow are:
\begin{enumerate}[label=(\alph*)]
\item consider only quasi-compact open subsets and closed subsets defined by ideals of finite type;
\item never take a scheme-theoretic closure, since that would take us outside (a); instead, take only a closed subset containing the given part and sufficiently small for the later constructions;
\item do not speak of a dense quasi-projective open subset; cf. (b). Remember that, to construct such open subsets, one could start with an affine cover $U_i$ and take $U_i'=U_i-\bigcup_{j<i}U_j$. Translate the arguments in terms of the original $U_i$, and apply (a)+(b).
\end{enumerate}

On the other hand, I have no idea in the case of algebraic spaces.

\section*{Acknowledgments}

The original notes were handwritten. They were typed by V. Maillot.

\begin{thebibliography}{9}
\bibitem[1]{conrad}
B. Conrad, Deligne's notes on Nagata compactifications, J. Ramanujan Math. Soc. \textbf{22} (2007), 205--257.

\bibitem[2]{nagata1}
M. Nagata, Imbedding of an abstract variety in a complete variety, J. Math. Kyoto Univ. \textbf{2} (1962), 1--10.

\bibitem[3]{nagata2}
M. Nagata, A generalization of the imbedding problem of an abstract variety in a complete variety, J. Math. Kyoto Univ. \textbf{3} (1963), 89--102.
\end{thebibliography}

\bigskip
\noindent Institute for Advanced Study, Princeton, New Jersey 08540, USA;\
deligne@math.ias.edu

\end{document}
